\documentclass[11pt,a4paper]{article}
\usepackage[T1]{fontenc}
\usepackage[utf8]{inputenc}
\usepackage{lmodern}
\usepackage{amsmath,amssymb,amsthm}
\usepackage{booktabs}
\usepackage{xcolor}
\usepackage[margin=27mm]{geometry}
\usepackage[colorlinks=true,urlcolor=blue,citecolor=blue,linkcolor=blue]{hyperref}
\newtheorem{theorem}{Theorem}
\newtheorem{proposition}[theorem]{Proposition}
\newtheorem{lemma}[theorem]{Lemma}
\newtheorem{corollary}[theorem]{Corollary}
\theoremstyle{definition}
\newtheorem{definition}[theorem]{Definition}
\newtheorem{example}[theorem]{Example}
\newtheorem{remark}[theorem]{Remark}
\theoremstyle{plain}
\DeclareMathOperator{\Tr}{Tr}
\DeclareMathOperator{\ran}{ran}
\DeclareMathOperator{\rank}{rank}
\newcommand{\ket}[1]{\lvert#1\rangle}
\newcommand{\bra}[1]{\langle#1\rvert}
\newcommand{\braket}[2]{\langle#1\vert#2\rangle}
\newcommand{\norm}[1]{\lVert#1\rVert}
\newcommand{\one}{\mathbf{1}}
\newcommand{\etac}{\eta_{\mathrm{crit}}}
\setlength{\emergencystretch}{2em}
\setlength{\parindent}{0pt}
\setlength{\parskip}{0.3em plus 0.1em minus 0.1em}
\title{How much leakage makes the Margenau--Hill state\\
of a post-selected system positive}
\author{Fable \& Astra\\[0.35em]
\normalsize Research led by Jean-Philippe Bruneton\\[-0.1em]
\small Universit\'e Paris Cit\'e, France}
\date{19 November 2026}

\begin{document}
\maketitle
\begin{abstract}
A quantum system observed between a preparation and a later post-selection
is described by two objects: the state $\rho$ inferred from the past and the
effect $E$ of the future outcome. A natural candidate for its state at the
intermediate time is the Margenau--Hill combination $(E\rho+\rho E)/2$,
normalized by the outcome probability, but this operator can have negative
eigenvalues. Kam and Wong showed that this happens whenever the post-selection
forbids a subspace that is not invariant under $\rho$. Real post-selections
are imperfect: a forbidden subspace leaks with some small probability $\eta$.
We determine exactly how much leakage restores positivity. For the effect
$E=P+\eta(\one-P)$, where $P$ is an orthogonal projector, write
$\rho$ in blocks $A$, $C$, $D$ relative to $P$ and let
$\kappa=\norm{A^{+1/2}CD^{+1/2}}$ be the largest normalized coherence between
the two blocks, with Moore--Penrose inverses. Then $\kappa\le1$, and the
Margenau--Hill operator is positive semidefinite if and only if
\[
\eta\ \ge\ \etac=\frac{1-\sqrt{1-\kappa^2}}{1+\sqrt{1-\kappa^2}}
=\frac{\kappa^2}{4}+\frac{\kappa^4}{8}+O(\kappa^6),
\]
in any finite dimension. Below the threshold we give an explicit vector on
which the operator is negative. The threshold is below one exactly when
$\rank\rho=\rank A+\rank D$, so for a pure state with weight in both blocks no
leakage short of $\eta=1$ suffices. For a qubit with Bloch vector
$(4/5,0,0)$ and $P=\ket0\bra0$, the threshold is exactly $1/4$.
\end{abstract}

\section{Introduction}
A system that is prepared, left to evolve, and then post-selected on a later
measurement outcome carries information from both ends of the time interval.
Quantum state smoothing \cite{Guevara2015} asks for the best estimate of the state at an intermediate
time given both a past and a future measurement record. In the simplest
description the past is summarized by a density matrix $\rho$, the filtered
state, and the future by an effect $0\le E\le\one$, the retrofiltered effect.
Combining the two into a single operator that behaves like a state is not
straightforward, because $\rho$ and $E$ need not commute.

Kam and Wong \cite{KamWong2026} study the family of orderings
$\Phi_s(E,\rho)=\tfrac12\bigl(\rho^sE\rho^{1-s}+\rho^{1-s}E\rho^s\bigr)$,
$s\in[0,1]$. All members have trace $\Tr(E\rho)$, the probability of the
post-selected outcome. The symmetric member $\rho^{1/2}E\rho^{1/2}$ is always
positive. The end point
\[
\Phi_0(E,\rho)=\tfrac12(E\rho+\rho E)
\]
is the Margenau--Hill ordering \cite{KamWong2026}, and it can fail to be
positive. Their Theorem~2 states that if $\Phi_0(E,\rho)\ge0$ then $\ker E$ is
invariant under $\rho$; equivalently, a forbidden subspace $\ker E$ that is
not $\rho$-invariant gives $\Phi_0$ a negative eigenvalue.

A post-selection that forbids a subspace exactly is an idealization. In
practice the forbidden subspace leaks: states in it pass the post-selection
with some small probability $\eta$. The simplest model is the effect
\[
E_\eta=P+\eta(\one-P),\qquad 0\le\eta\le1,
\]
where $P$ is an orthogonal projector; for $\eta=0$ the forbidden subspace is
$\ran(\one-P)$, and for $\eta=1$ there is no post-selection at all and
$\Phi_0=\rho$. This note answers the question: what is the smallest leakage
$\eta$ for which the Margenau--Hill operator $\Phi_0(E_\eta,\rho)$ is positive?
Since dividing by $\Tr(E_\eta\rho)>0$ does not affect positivity, the same
threshold governs the normalized Margenau--Hill state.

The answer depends on $\rho$ through a single number. Writing $\rho$ in blocks
$A=P\rho P$, $C=P\rho(\one-P)$, $D=(\one-P)\rho(\one-P)$, the quantity
$\kappa=\norm{A^{+1/2}CD^{+1/2}}$ measures the largest coherence between the
kept and the leaking subspaces, relative to their populations. We show that
$\kappa\le1$ and that $\Phi_0(E_\eta,\rho)\ge0$ exactly when
$4\eta/(1+\eta)^2\ge\kappa^2$, that is, when $\eta$ exceeds
$\etac=(1-\sqrt{1-\kappa^2})/(1+\sqrt{1-\kappa^2})$. Weak coherence needs
leakage of order $\kappa^2/4$; maximal coherence $\kappa=1$, which occurs for
every pure state with weight in both subspaces, cannot be compensated by any
leakage below one.

A single qubit already displays the whole phenomenon. Take
$\rho=\tfrac12(\one+\tfrac45X)$ and $P=\ket0\bra0$. The Margenau--Hill
operator is
$\Phi_0=\tfrac12\left(\begin{smallmatrix}1&\frac{2}{5}(1+\eta)\\[0.2em]
\frac25(1+\eta)&\eta\end{smallmatrix}\right)$, with determinant
$-(4\eta-1)(\eta-4)/100$. It is negative for $\eta<1/4$ and positive for
$1/4\le\eta\le1$, and indeed $\kappa=4/5$ gives $\etac=1/4$. The general
result replaces the off-diagonal entry by the matrix $A^{+1/2}CD^{+1/2}$ and its
modulus by the operator norm.

\section{Related work}\label{sec:related}
Guevara and Wiseman
\cite{Guevara2015} formulated quantum state smoothing, the estimation of a
quantum state from past and future records, which is the setting in which
operators built from a filtered state $\rho$ and a retrofiltered effect $E$
arise.

The question answered here is a quantitative refinement of the work of Kam and
Wong \cite{KamWong2026}. Their Theorem~1 shows, for full-rank $\rho$ and an
effect with nontrivial kernel, that an open interval of orderings $\Phi_s$
around $s=1/2$ remains positive if and only if $\ker E$ is $\rho$-invariant, and
their Theorem~2 shows that the Margenau--Hill member $\Phi_0$ cannot be positive
when $\ker E$ is not $\rho$-invariant. Both statements are qualitative and
concern effects with a kernel. A leaking effect $E_\eta$ with $\eta>0$ has
trivial kernel, so their theorems do not apply to it, and the present note
measures how far the obstruction persists. At $\eta=0$ our criterion reduces
to $C=0$, which is precisely the invariance of $\ker E_0=\ran(\one-P)$ under
$\rho$; for this family of effects it thus recovers their Theorem~2 together
with its converse (Remark~\ref{rem:eta0}).

The proofs use the Schur-complement characterization of positive semidefinite
block matrices with a possibly singular diagonal block. It is classical; we
include the short proof in Lemma~\ref{lem:schur} so that the note is
self-contained.

\section{Setting}\label{sec:setting}
Let $\mathcal H$ be a finite-dimensional Hilbert space, $\rho\ge0$ a density
matrix on $\mathcal H$, and $P$ an orthogonal projector with $0\ne P\ne\one$.
Relative to the decomposition $\mathcal H=\ran P\oplus\ran(\one-P)$ write
\[
\rho=\begin{pmatrix}A&C\\ C^\dagger&D\end{pmatrix},\qquad
E_\eta=\begin{pmatrix}\one&0\\0&\eta\one\end{pmatrix}.
\]
A direct multiplication gives
\begin{equation}\label{eq:image}
\Phi(\eta):=\Phi_0(E_\eta,\rho)=\tfrac12(E_\eta\rho+\rho E_\eta)
=\begin{pmatrix}A&\beta C\\ \beta C^\dagger&\eta D\end{pmatrix},
\qquad\beta=\frac{1+\eta}{2}.
\end{equation}
For a positive semidefinite matrix $M$ we write $M^+$ for its Moore--Penrose
inverse, $M^{+1/2}=(M^{1/2})^+$, and $\Pi_M=MM^+$ for the orthogonal projector
onto $\ran M$. Define
\[
K=A^{+1/2}\,C\,D^{+1/2},\qquad \kappa=\norm{K},
\]
the operator norm. When $A$ and $D$ are invertible, $K=A^{-1/2}CD^{-1/2}$ and
$\kappa$ is the largest singular value of the coherence block after each
diagonal block has been normalized to the identity.

\begin{lemma}[Schur complement]\label{lem:schur}
Let $M=\left(\begin{smallmatrix}A&B\\B^\dagger&F\end{smallmatrix}\right)$ be
Hermitian with $A\ge0$. Then $M\ge0$ if and only if $\ran B\subseteq\ran A$
and $F-B^\dagger A^+B\ge0$. In that case
\begin{equation}\label{eq:congruence}
M=L^\dagger\begin{pmatrix}A&0\\0&F-B^\dagger A^+B\end{pmatrix}L,\qquad
L=\begin{pmatrix}\one&A^+B\\0&\one\end{pmatrix},
\end{equation}
and for every vector $v$ the vector $u=-A^+Bv$ satisfies
$\langle(u,v),M(u,v)\rangle=\langle v,(F-B^\dagger A^+B)v\rangle$.
\end{lemma}

\begin{proof}
Suppose $M\ge0$ and let $x\in\ker A$. For a real $t$ and any $v$,
$\langle(tx,v),M(tx,v)\rangle=2t\,\mathrm{Re}\langle x,Bv\rangle+\langle v,Fv\rangle\ge0$,
which forces $\langle x,Bv\rangle=0$. Hence $\ker A\subseteq\ker B^\dagger$, that
is, $\ran B\subseteq\ran A$. Whenever this inclusion holds, $AA^+B=B$, and
expanding $L^\dagger\,\mathrm{diag}(A,F-B^\dagger A^+B)\,L$ gives back $M$, which
is \eqref{eq:congruence}. Since $L$ is invertible, $M\ge0$ is then equivalent
to $F-B^\dagger A^+B\ge0$. The last identity follows by expanding
$\langle u,Au\rangle+2\,\mathrm{Re}\langle u,Bv\rangle+\langle v,Fv\rangle$ with
$u=-A^+Bv$ and $A^+AA^+=A^+$.
\end{proof}

\begin{lemma}[Factorization of the coherence]\label{lem:factor}
If $\rho\ge0$ then $\ran C\subseteq\ran A$ and $\ran C^\dagger\subseteq\ran D$.
Consequently $C=A^{1/2}KD^{1/2}$, $C^\dagger A^+C=D^{1/2}K^\dagger KD^{1/2}$, and
$K^\dagger K=\Pi_DK^\dagger K\Pi_D$.
\end{lemma}

\begin{proof}
The first inclusion is Lemma~\ref{lem:schur} applied to $\rho$; the second is
the same lemma applied to $\rho$ with the two blocks interchanged. The
inclusions give $\Pi_AC=C=C\Pi_D$. Since $A^{1/2}A^{+1/2}=\Pi_A$ and
$D^{+1/2}D^{1/2}=\Pi_D$, we get $A^{1/2}KD^{1/2}=\Pi_AC\Pi_D=C$. Then
$C^\dagger A^+C=D^{1/2}K^\dagger A^{1/2}A^+A^{1/2}KD^{1/2}=D^{1/2}K^\dagger\Pi_AKD^{1/2}$,
and $\Pi_AK=K$ because $\ran A^{+1/2}=\ran A$. Finally $K\Pi_D=K$ because
$D^{+1/2}\Pi_D=D^{+1/2}$, which gives the last identity.
\end{proof}

\section{The threshold}\label{sec:main}
\begin{theorem}\label{thm:main}
Let $\rho\ge0$, $P$ and $\kappa$ be as above. Then $\kappa\le1$, and for
$0\le\eta\le1$ the Margenau--Hill operator $\Phi(\eta)$ of \eqref{eq:image} is
positive semidefinite if and only if
\begin{equation}\label{eq:criterion}
\frac{4\eta}{(1+\eta)^2}\ \ge\ \kappa^2,
\end{equation}
that is, if and only if $\eta\ge\etac$ with
\[
\etac=\frac{1-\sqrt{1-\kappa^2}}{1+\sqrt{1-\kappa^2}}.
\]
\end{theorem}

\begin{proof}
If $D=0$ then $C=0$ by Lemma~\ref{lem:factor}, so $\kappa=0$, $\etac=0$, and
$\Phi(\eta)=\mathrm{diag}(A,0)\ge0$ for every $\eta$; the statement holds.
Assume $D\ne0$. Apply Lemma~\ref{lem:schur} to $\Phi(\eta)$ with $B=\beta C$ and
$F=\eta D$. The range condition $\ran(\beta C)\subseteq\ran A$ holds by
Lemma~\ref{lem:factor}, so $\Phi(\eta)\ge0$ if and only if
\[
S_\eta:=\eta D-\beta^2C^\dagger A^+C
=D^{1/2}\bigl(\eta\,\Pi_D-\beta^2K^\dagger K\bigr)D^{1/2}\ \ge\ 0,
\]
where the factorization uses Lemma~\ref{lem:factor} and $D^{1/2}\Pi_DD^{1/2}=D$.
The operator $\eta\Pi_D-\beta^2K^\dagger K$ is supported on $\ran D$, on which
$D^{1/2}$ is invertible, so $S_\eta\ge0$ if and only if
$\eta\Pi_D\ge\beta^2K^\dagger K$. As $K^\dagger K$ is supported on $\ran D$ and
its largest eigenvalue is $\kappa^2$, this holds if and only if
$\eta\ge\beta^2\kappa^2$, which is \eqref{eq:criterion}.

For $\eta=1$ we have $\Phi(1)=\rho\ge0$, and \eqref{eq:criterion} gives
$\kappa\le1$. The function $f(\eta)=4\eta/(1+\eta)^2$ is continuous and strictly
increasing on $[0,1]$, since $f'(\eta)=4(1-\eta)/(1+\eta)^3>0$ there, with
$f(0)=0$ and $f(1)=1$. The set of $\eta\in[0,1]$ satisfying
\eqref{eq:criterion} is therefore an interval $[\eta_*,1]$, where $\eta_*$ is
the unique solution of $f(\eta)=\kappa^2$ in $[0,1]$. Writing
$\sigma=\sqrt{1-\kappa^2}$, the equation $\kappa^2(1+\eta)^2=4\eta$ has the two
roots $(1-\sigma)/(1+\sigma)$ and $(1+\sigma)/(1-\sigma)$, whose product is
$1$; the first lies in $[0,1]$, so $\eta_*=\etac$. (For $\kappa=0$ the
criterion holds for every $\eta$ and $\etac=0$.)
\end{proof}

Three features of the threshold are immediate. It depends on $\rho$ and $P$
only through $\kappa$, and not on the dimension. It increases from $0$ at
$\kappa=0$ to $1$ at $\kappa=1$, with expansion
\[
\etac=\frac{\kappa^2}{4}+\frac{\kappa^4}{8}+\frac{5\kappa^6}{64}+O(\kappa^8),
\]
so weak coherence is compensated by leakage of order $\kappa^2/4$. And the
positive set is the whole interval $[\etac,1]$: more leakage never destroys
positivity once it has been reached.

\begin{remark}\label{rem:eta0}
At $\eta=0$ the criterion reads $\kappa=0$, which by Lemma~\ref{lem:factor} is
equivalent to $C=0$, that is, to the invariance of $\ker E_0=\ran(\one-P)$ under
$\rho$. For this family of effects the theorem therefore contains Theorem~2 of
\cite{KamWong2026} and its converse: $\Phi_0(E_0,\rho)\ge0$ if and only if
$\ker E_0$ is $\rho$-invariant. The note does not require $\rho$ to have full
rank.
\end{remark}

\section{The negative direction below threshold}\label{sec:negative}
The proof of Theorem~\ref{thm:main} also produces an explicit vector on which
$\Phi(\eta)$ is negative when $\eta<\etac$. Assume $\kappa>0$, and let $y$ be a
unit right singular vector of $K$ for its largest singular value, so that
$\norm{Ky}=\kappa$. Since $y\in\ran K^\dagger\subseteq\ran D$, the vector
$v=D^{+1/2}y$ satisfies $\langle v,Dv\rangle=\norm{\Pi_Dy}^2=1$. Put
\[
u=-\beta A^+Cv,\qquad w=(u,v)\in\ran P\oplus\ran(\one-P).
\]

\begin{proposition}\label{prop:witness}
For every $\eta\in[0,1]$,
\[
\langle w,\Phi(\eta)w\rangle=\eta-\frac{(1+\eta)^2}{4}\,\kappa^2
=\frac{(1+\eta)^2}{4}\bigl(f(\eta)-\kappa^2\bigr),
\]
which is negative exactly when $\eta<\etac$.
\end{proposition}

\begin{proof}
By the last statement of Lemma~\ref{lem:schur} with $B=\beta C$ and $F=\eta D$,
$\langle w,\Phi(\eta)w\rangle=\langle v,S_\eta v\rangle$. With the
factorization of $S_\eta$,
$\langle v,S_\eta v\rangle=\langle y,\Pi_DD^{1/2}D^{+1/2}(\eta\Pi_D-\beta^2K^\dagger K)D^{+1/2}D^{1/2}\Pi_Dy\rangle
=\eta\norm{y}^2-\beta^2\norm{Ky}^2=\eta-\beta^2\kappa^2$, because
$D^{1/2}D^{+1/2}=\Pi_D$ and $\Pi_Dy=y$. The sign statement is
Theorem~\ref{thm:main}.
\end{proof}

The witness has a transparent structure. Its component $v$ in the leaking
subspace is the direction of maximal normalized coherence, and its component
$u$ in the kept subspace is the best linear response to $v$, which is how the
Schur complement isolates the off-diagonal contribution. The depth of the
negativity, $\eta-\beta^2\kappa^2$, is largest at $\eta=0$, where it equals
$-\kappa^2/4$ for this normalization of $w$.

\section{When the threshold is below one}\label{sec:rank}
\begin{proposition}\label{prop:rank}
$\kappa<1$ if and only if $\rank\rho=\rank A+\rank D$. In particular $\kappa<1$
whenever $\rho$ has full rank, and $\kappa=1$ for every pure state
$\rho=\ket\psi\bra\psi$ with $P\ket\psi\ne0$ and $(\one-P)\ket\psi\ne0$.
\end{proposition}

\begin{proof}
If $D=0$ then $\kappa=0$ and $\rank\rho=\rank A$, since $C=0$. Otherwise,
applying \eqref{eq:congruence} to $\rho$ itself shows that $\rho$ is congruent
to $\mathrm{diag}(A,S_1)$ with $S_1=D^{1/2}(\Pi_D-K^\dagger K)D^{1/2}$, so
$\rank\rho=\rank A+\rank S_1$. Since $0\le K^\dagger K\le\Pi_D$ and $D^{1/2}$ is
invertible on $\ran D$, $\rank S_1=\rank D$ exactly when $\Pi_D-K^\dagger K$ is
invertible on $\ran D$, that is, when $\kappa<1$. For full-rank $\rho$ the
blocks $A$ and $D$ are invertible and the ranks add. For the pure state, $A$
and $D$ have rank one while $\rho$ has rank one, so $\kappa=1$.
\end{proof}

The second case has a direct physical reading: if the system is in a pure
state with amplitude in both the kept and the leaking subspace, the
Margenau--Hill operator of any genuinely selective leaky post-selection,
$\eta<1$, has a negative eigenvalue. A leakage below one can restore
positivity only when the ranks of the two diagonal blocks add up to the rank
of $\rho$, which is the case for every full-rank state.

\section{Qubit example}\label{sec:qubit}
\begin{example}\label{ex:qubit}
Let $\rho=\tfrac12(\one+xX+yY+zZ)$ with $x^2+y^2+z^2\le1$ and $\lvert z\rvert<1$,
and $P=\ket0\bra0$. Then $A=(1+z)/2$, $D=(1-z)/2$ and $C=(x-iy)/2$, so
\[
\kappa^2=\frac{x^2+y^2}{1-z^2},\qquad
\det\Phi(\eta)=\frac{\eta(1-z^2)-\beta^2(x^2+y^2)}{4}.
\]
The condition $\kappa<1$ says that the Bloch vector lies strictly inside the
sphere, in agreement with Proposition~\ref{prop:rank}. For $(x,y,z)=(4/5,0,0)$
one finds $\kappa=4/5$, $\sqrt{1-\kappa^2}=3/5$ and $\etac=(2/5)/(8/5)=1/4$; the
determinant equals $-(4\eta-1)(\eta-4)/100$, which vanishes at $\eta=1/4$ and at
the reciprocal root $\eta=4$. For $(x,y,z)=(3/5,0,0)$ one finds $\etac=1/9$.
In this example the explicit witness of Proposition~\ref{prop:witness} is
$v=(1-z)^{-1/2}\sqrt2$ and $u=-\beta(x-iy)\,v/(1+z)$, up to a phase.
\end{example}

\section{Verification}\label{sec:verification}
The results above are proved analytically. The public
\href{https://github.com/jpbruneton/open_quantum_problems/blob/main/research/entries/postselection-leakage-threshold/checks/verify_leakage_threshold.py}{Python checker}
tests them independently and writes its
\href{https://github.com/jpbruneton/open_quantum_problems/blob/main/research/entries/postselection-leakage-threshold/data/verification.json}{recorded output}
into the same entry. In exact symbolic arithmetic it confirms that $\etac$ is
the smaller root of $\kappa^2(1+\eta)^2=4\eta$, that the two roots are
reciprocal, the series expansion of $\etac$, and the determinant and
thresholds of Example~\ref{ex:qubit}. In floating-point arithmetic it draws
random density matrices of dimension $2$ to $6$ (full rank, rank deficient and
pure) with random projectors of every rank, $900$ instances in all, together
with $200$ instances in which both diagonal blocks are singular. For each it
computes $\kappa$ with Moore--Penrose inverses and checks that $\kappa\le1$,
that the smallest eigenvalue of $\Phi(\eta)$ is nonnegative on $[\etac,1]$,
that the witness of Proposition~\ref{prop:witness} takes the predicted value
and that $\Phi(\eta)$ is negative below $\etac$, that a bisection on the sign of
the smallest eigenvalue reproduces $\etac$, and that $\kappa<1$ exactly when
$\rank\rho=\rank A+\rank D$. The bisection agrees with the formula to
$3\times10^{-10}$ when $\kappa<1$; when $\kappa=1$ the smallest eigenvalue
vanishes to second order at $\eta=1$ and the agreement is $10^{-4}$, while the
witness confirms negativity for every tested $\eta<1$. The checker uses a fixed
seed and runs in a few seconds.

\section{Scope and limitations}\label{sec:scope}
Theorem~\ref{thm:main} and Propositions~\ref{prop:witness}
and~\ref{prop:rank} are proved analytically in any finite dimension, for
arbitrary density matrices including singular ones. The computations confirm
these statements numerically and verify the exact examples; they are not a
substitute for the proofs.

The result concerns the Margenau--Hill ordering $s=0$ and the isotropic
leakage model $E_\eta=P+\eta(\one-P)$. Two natural extensions are left open:
the positivity of the intermediate orderings $\Phi_s$ with $0<s<1/2$ under
leakage, which governs the interval of admissible orderings studied in
\cite{KamWong2026}, and anisotropic leakage, in which $\eta(\one-P)$ is
replaced by a general effect $F$ supported on $\ran(\one-P)$. In the latter
case the lower-right block of $\Phi_0$ becomes $(FD+DF)/2$ and the off-diagonal
block $C(\one+F)/2$, and the reduction to a single scalar $\kappa$ no longer
applies.

On bibliographic priority, the threshold formula was not found in the
documented search, which consisted of reading the work of Kam and Wong
\cite{KamWong2026} and of Guevara and Wiseman \cite{Guevara2015}. The result is
an elementary consequence of the Schur-complement criterion, and a related
statement may exist elsewhere; this note does not certify priority.

\begin{thebibliography}{9}
\setlength{\itemsep}{2pt}
\bibitem{Guevara2015}
I. Guevara and H. Wiseman,
\emph{Quantum state smoothing},
Phys. Rev. Lett. \textbf{115}, 180407 (2015);
\href{https://doi.org/10.1103/PhysRevLett.115.180407}{doi:10.1103/PhysRevLett.115.180407};
\href{https://arxiv.org/abs/1503.02799}{arXiv:1503.02799}.

\bibitem{KamWong2026}
C.-F. Kam and K.-W. Wong,
\emph{Forbidden subspaces in quantum state smoothing},
arXiv preprint (2026);
\href{https://arxiv.org/abs/2609.17307}{arXiv:2609.17307}.
\end{thebibliography}
\end{document}
